\chapter{Torsja Whiteheada i twierdzenie o \texorpdfstring{$s$}{s}-kobordyzmie}
\label{ch:s-kobordyzm}
\index{s-kobordyzm@$s$-kobordyzm}

W Rozdziale~\ref{ch:hcobordyzm} zakładaliśmy prostą spójność.
Pozwalała ona zastępować zbiory przecięć jedną liczbą całkowitą.
Jeśli $\pi_1(W)$ jest nietrywialna, punkt przecięcia niesie jeszcze
informację o drodze, którą do niego dotarliśmy. Dwa punkty
z przeciwnymi znakami mogą więc \emph{nie} być parą nadającą się
do usunięcia. Torsja Whiteheada zapisuje algebraiczny ślad
tej trudności po wszystkich dopuszczalnych ruchach uchwytów.
Litera $s$ pochodzi od angielskiego \emph{simple homotopy}.

Przez cały rozdział $W^m$ jest spójnym zwartym gładkim kobordyzmem
$(W;N_-,N_+)$ o niepustych, spójnych, zamkniętych brzegach,
$\pi=\pi_1(W)$, a $R=\Z[\pi]$. W twierdzeniu o
$s$-kobordyzmie przyjmiemy $m\geq6$. W podstawowej części
nie zakładamy orientowalności $W$.

\section{Po co nam pierścień grupowy i nakrycie uniwersalne?}
\label{sec:pi-lancuchy-20}

\begin{definicja}[Pierścień grupowy]
\index{pierscien grupowy@Pierścień grupowy}
\label{def:pierscien-grupowy-20}
Element $R=\Z[\pi]$ jest skończoną sumą formalną
$\sum_{g\in\pi} n_g g$, gdzie $n_g\in\Z$ i niemal wszystkie
$n_g$ są zerowe. Dodajemy współczynniki przy tym samym $g$,
a mnożenie określa rozdzielność oraz reguła
$g\cdot h=gh$ w grupie $\pi$. Jednostką pierścienia jest
element neutralny $e\in\pi$. Jeśli $\pi$ nie jest przemienna,
to na ogół również $R$ nie jest przemienny.
\end{definicja}

\begin{przyklad}[Okrąg i grupa cykliczna]
Dla $\pi=\Z=\langle t\rangle$ otrzymujemy
$\Z[\pi]=\Z[t,t^{-1}]$: wolno używać skończonych sum
potęg dodatnich i ujemnych $t$. Jeśli $\pi=C_5=\langle
t\mid t^5=e\rangle$, to
$\Z[C_5]\cong\Z[t]/(t^5-1)$. Wyrażenia $2-t+t^4$ oraz
$2-t+t^{-1}$ oznaczają w drugim pierścieniu ten sam element.
\end{przyklad}

Niech $p:\widetilde W\to W$ będzie nakryciem uniwersalnym,
a $\widetilde N_-=p^{-1}(N_-)$. Wybierając konwencję
$\widetilde x\cdot g=g^{-1}\widetilde x$ dla działania
przekształceń nakrycia, grupy łańcuchów komórkowych pary
$(\widetilde W,\widetilde N_-)$ stają się \emph{prawymi}
modułami nad $R$. Z każdego uchwytu wybieramy jedno
podniesienie jego rdzenia. Te wybrane podniesienia, po jednym
na uchwyt, tworzą bazę nad $R$. Wszystkie ich translacje przez
$\pi$ tworzą natomiast bazę nad $\Z$. Zmiana podniesienia mnoży
odpowiedni wektor bazowy przez element $g\in\pi$.
Przy ustalonej orientacji rdzenia zmiana orientacji mnoży go
przez $-1$. Zatem nieoznaczone naturalne wybory bazy
mają dokładnie postać $\pm g$.

\begin{stwierdzenie}[Względny kompleks uchwytów nad $R$]
\label{prop:kompleks-pi-20}
Dla rozkładu $W$ na uchwyty względem $N_-$ istnieje
skończony kompleks wolnych prawych modułów
\[
 0\longrightarrow C_m\xrightarrow{\partial_m}\cdots
 \xrightarrow{\partial_2}C_1\xrightarrow{\partial_1}C_0
 \longrightarrow0,
 \qquad C_k\cong R^{c_k},
\]
gdzie $c_k$ jest liczbą $k$-uchwytów. Oblicza on homologie
$H_*(\widetilde W,\widetilde N_-;\Z)$, a wybrane podniesienia
i orientacje rdzeni stanowią jego bazy.
\end{stwierdzenie}
\begin{proof}
Filtrujmy $W$ podpoziomicami po kolejnych indeksach uchwytów
i podnieśmy tę filtrację do $\widetilde W$. Każde podniesienie
$k$-uchwytu daje generator grupy względnej filtracji w stopniu
$k$, a podniesienia jednego uchwytu tworzą orbitę działania
$\pi$. Jest to jeden egzemplarz wolnego modułu $R$.
Różniczka jest homomorfizmem łączącym dla trzech kolejnych
podpoziomic, jak w dowodzie
Twierdzenia~\ref{thm:kompleks-uchwytow-18}. Jej kwadrat
jest zerem z dokładności ciągów par, a homologii otrzymanego
kompleksu odpowiada homologia pary nakryć przez argument
filtracyjny z tego twierdzenia. Jedyna zmiana względem
kompleksu nad $\Z$ polega na zachowaniu oznaczeń wszystkich
podniesionych uchwytów.
\end{proof}

\begin{uwaga}[Acykliczność nad $R$]
Jeżeli $N_-\hookrightarrow W$ jest równoważnością
homotopijną, to jej podniesienie
$\widetilde N_-\hookrightarrow\widetilde W$ także jest
równoważnością homotopijną. Dlatego kompleks z
poprzedniego stwierdzenia jest acykliczny. To silniejszy
fakt niż samo $H_*(W,N_-;\Z)=0$: w tym drugim rachunku
wszystkie $g$ zostają zastąpione przez $1$ za pomocą
\emph{augmentacji} $\epsilon:R\to\Z$, $\sum n_g g\mapsto
\sum n_g$.
\end{uwaga}

Ustalmy poziomicę między uchwytami indeksów $k$ i $k+1$,
jak w Rozdziale~\ref{ch:rachunek-uchwytow}. Niech $A_j$
oznacza sferę przyczepienia $(k+1)$-uchwytu, a $B_i$
sferę pasa $k$-uchwytu. Po wyborze podniesień i dróg
do punktu bazowego każdemu przecięciu $x\in A_j\cap B_i$
odpowiada znak $\varepsilon_x\in\{\pm1\}$ i element
$g_x\in\pi$: wybrane podniesienie $A_j$ spotyka w punkcie
nad $x$ prawą translację $\widetilde B_i\cdot g_x
=g_x^{-1}\widetilde B_i$ wybranej sfery pasa.
Orientację podniesionej sfery przyczepienia i orientację normalną
sfery pasa bierzemy z orientacji odpowiednich rdzeni.
Przenosimy je razem z translacjami; nie wybieramy globalnej
orientacji $W$, która może nie istnieć.
Przy naszych bazach
\begin{equation}\label{eq:macierz-pi-20}
 \partial_{k+1}(e_j)=\sum_i f_i a_{ij},
 \qquad a_{ij}=\sum_{x\in A_j\cap B_i}\varepsilon_x g_x
 \ \in R.
\end{equation}
Prawa strona jest \emph{sumą w pierścieniu grupowym}.
Dopiero jej augmentacja daje podpisaną liczbę przecięć
z~\eqref{eq:macierz-uchwyty-18}.

\begin{figure}[htbp]
\centering
\begin{tikzpicture}[>=Latex,font=\small,
 every node/.style={align=center,rounded corners=3pt,inner sep=5pt}]
 \node[draw=manifoldblue!75!black,fill=manifoldblue!9]
 (a) at (-3.05,0) {punkty $x,y\in A_j\cap B_i$\\
                    znaki $+,-$};
 \node[draw=manifoldgold!80!black,fill=manifoldgold!12]
 (b) at (.08,0) {etykiety dróg\\$g_x,\ g_y\in\pi$};
 \node[draw=manifoldgreen!75!black,fill=manifoldgreen!10]
 (c) at (3.30,0) {$g_x=g_y$\\pętla ściągalna};
 \node[draw=manifoldred!75!black,fill=manifoldred!9]
 (d) at (.08,-1.25) {$g_x\ne g_y$\\pętla może przeszkadzać};
 \draw[->,thick] (a)--(b);
 \draw[->,thick,manifoldgreen!70!black] (b)--(c);
 \draw[->,thick,manifoldred!75!black] (b)--(d);
\end{tikzpicture}
\caption{Znaki $+,-$ dają zerową zwykłą liczbę przecięcia.
Iloczyn $g_xg_y^{-1}$ opisuje pętlę powstałą z łuków na obu
sferach, z dokładnością do sprzężenia i wyboru kierunku.
Równe etykiety pozwalają usunąć parę przy założeniach
wymiarowych i geometrycznych Lematu~\ref{lem:whitney-pi-20}.}
\label{fig:etykiety-whitneya-20}
\end{figure}

\begin{przyklad}[Dlaczego augmentacja traci dane?]
Dwa punkty o znakach $+,-$ i etykietach $e,t\in\pi$
dają w~\eqref{eq:macierz-pi-20} wyraz $1-t$.
Po augmentacji jest to $1-1=0$, lecz w $\Z[\pi]$
może być niezerowy. Łuk Whitneya łączący te punkty
ma wówczas pętlę o klasie $t$ (lub $t^{-1}$, zależnie
od kierunku), zatem nie uzyskaliśmy dysku przez sam
argument o przeciwnych znakach.
\end{przyklad}

\section{Grupa \texorpdfstring{$K_1$}{K1} i grupa Whiteheada}
\label{sec:k1-wh-20}

Nad $\Z$ redukcja za pomocą algorytmu Euklidesa używa dodawania wierszy,
permutacji oraz mnożenia wiersza przez $-1$. Nad
$R=\Z[\pi]$ nie ma ogólnie wyznacznika, ale nadal są
odwracalne macierze i ruchy odpowiadające przesunięciom
uchwytów. Porównujemy więc macierze po dopuszczeniu
bloków jednostkowych.

\begin{definicja}[$K_1$ i $\operatorname{Wh}(\pi)$]
\label{def:k1-wh-20}
\index{K1@$K_1$}
\index{grupa Whiteheada@Grupa Whiteheada}
Niech $GL_r(R)$ będzie grupą odwracalnych macierzy
$r\times r$ nad $R$. Włączamy $GL_r(R)$ do $GL_{r+1}(R)$
przez $A\mapsto\operatorname{diag}(A,1)$ i oznaczamy
sumę rosnącą tych grup przez $GL_\infty(R)$.
Każdy jej element ma więc reprezentanta pewnego skończonego rozmiaru.
Macierze elementarne
to $E_{ij}(a)=I+a e_{ij}$ dla $i\ne j$ i $a\in R$;
$E_\infty(R)$ jest generowaną przez nie podgrupą.
Definiujemy
\[
 K_1(R)=GL_\infty(R)/E_\infty(R),\qquad
 \operatorname{Wh}(\pi)
 =K_1(\Z[\pi])\big/\langle[\pm g]:g\in\pi\rangle.
\]
Grupę zapisujemy addytywnie: klasę macierzy $A$
oznaczamy $[A]$, a jej element zerowy $[I]$.
\end{definicja}

\begin{stwierdzenie}[Dlaczego iloraz jest grupą abelową?]
\label{prop:k1-abelowa-20}
Podgrupa $E_\infty(R)$ jest normalna w $GL_\infty(R)$,
a iloraz $K_1(R)$ jest abelowy. W szczególności
$[AB]=[A]+[B]$, $[A\oplus B]=[A]+[B]$ i
$[A^{-1}]=-[A]$.
\end{stwierdzenie}
\begin{proof}
Dla trzech różnych indeksów $i,j,k$ bezpośrednie
mnożenie macierzy daje
\[
 [E_{ik}(a),E_{kj}(b)]=E_{ij}(ab).
\]
Po powiększeniu rozmiaru każda macierz elementarna
jest więc komutatorem. Aby sprawdzić drugie
zawieranie, dla dowolnej odwracalnej macierzy
$A$ rozkładamy $\operatorname{diag}(A,A^{-1})$
na macierze elementarne. Istotnie, pomnożenie
\[
 \begin{pmatrix}I&A\\0&I\end{pmatrix}
 \begin{pmatrix}I&0\\-A^{-1}&I\end{pmatrix}
 \begin{pmatrix}I&A\\0&I\end{pmatrix}
 \begin{pmatrix}0&-I\\I&0\end{pmatrix}
 =\begin{pmatrix}A&0\\0&A^{-1}\end{pmatrix}
\]
daje żądaną macierz. Ostatni czynnik ma rozkład
\[
 \begin{pmatrix}0&-I\\I&0\end{pmatrix}
 =\begin{pmatrix}I&-I\\0&I\end{pmatrix}
  \begin{pmatrix}I&0\\I&I\end{pmatrix}
  \begin{pmatrix}I&-I\\0&I\end{pmatrix};
\]
każdy blokowy czynnik jednostkowo trójkątny
jest iloczynem zwykłych macierzy elementarnych.
Jeśli $G$ jest odwracalna, a $E$ elementarna,
to
\[
 \operatorname{diag}(GEG^{-1},I)
 =\operatorname{diag}(G,G^{-1})
  \operatorname{diag}(E,I)
  \operatorname{diag}(G^{-1},G)
 \in E_\infty(R).
\]
To dowodzi normalności $E_\infty(R)$.
Możemy więc utworzyć iloraz. Poprzedni rozkład
pokazuje, że $A\oplus I$ i $I\oplus A$ mają
w nim tę samą klasę. Bloki $A\oplus I$ i $I\oplus B$
komutują jako macierze, toteż ich klasy, czyli
klasy $A$ i $B$, komutują. Iloraz jest abelowy,
więc zawiera wszystkie komutatory w jądrze,
czyli $[GL_\infty(R),GL_\infty(R)]\subseteq E_\infty(R)$.
Razem z pierwszym zawieraniem dostajemy równość.
Dla bezpośredniej sumy używamy faktoryzacji
$A\oplus B=(A\oplus I)(I\oplus B)$.
\end{proof}

Macierz $E_{ij}(a)$ oznacza dodanie wielokrotności
jednego generatora do innego. Blok $I$ odpowiada
dopisaniu pary, która się znosi. Mnożnik $\pm g$
wyraża zmianę orientacji i podniesienia pojedynczej
komórki. Tych zmian nie chcemy uznawać za nową
przeszkodę. Dlatego $\operatorname{Wh}(\pi)$ jest
właściwym ilorazem $K_1(R)$ dla geometrii uchwytów.

\begin{stwierdzenie}[Przypadek prosto spójny]
\label{prop:wh-trywialna-20}
$\operatorname{Wh}(\{e\})=0$.
\end{stwierdzenie}
\begin{proof}
Wtedy $R=\Z$. Macierz całkowita odwracalna ma
wyznacznik $\pm1$. Algorytm Euklidesa w pierwszej
kolumnie, a potem indukcja po rozmiarze sprowadzają
ją dodawaniem całkowitych wielokrotności wierszy,
przestawianiem i zmianą znaków do $I$. Zamiana
wierszy jest iloczynem macierzy elementarnych oraz
jednego czynnika $-1$; po przejściu do
$\operatorname{Wh}(\{e\})$ ten czynnik także jest
zerowy. Zatem klasa każdej macierzy znika.
\end{proof}

\begin{przyklad}[Konkretny niezerowy element dla $C_5$]
\label{prz:wh-c5-20}
W $R=\Z[C_5]$ niech $t^5=1$ i
\[
 u=1-t-t^{-1},\qquad v=1-t^2-t^{-2}.
\]
Ponieważ $t^3=t^{-2}$ i $t^{-3}=t^2$, mamy
\[
 (t+t^{-1})(t^2+t^{-2})=t^2+t^{-2}+t+t^{-1}.
\]
Po rozwinięciu $(1-t-t^{-1})(1-t^2-t^{-2})$ wszystkie
wyrazy poza $1$ się znoszą. Zatem $uv=1$,
więc $(u)$ jest macierzą odwracalną rozmiaru $1$.
Nie jest to jedna z oczywistych jednostek $\pm t^j$.
Aby dowieść mocniejszego faktu $[(u)]\ne0$ w
$\operatorname{Wh}(C_5)$, wybierzmy
$\zeta=e^{2\pi i/5}$ i zastosujmy homomorfizm
$R\to\C$ wysyłający $t$ na $\zeta$.
Po przejściu do macierzy zespolonych ich wyznacznik
ma własność multiplikatywną. Macierze elementarne
mają wyznacznik $1$, bloki $I$ również, zaś
$|\det(\pm\zeta^j)|=1$. Tymczasem
\[
 |u(\zeta)|=|1-\zeta-\zeta^{-1}|
 =1-2\cos(2\pi/5)=\frac{3-\sqrt5}{2}\ne1.
\]
Moduł wyznacznika rozróżnia więc klasę $(u)$
od zera w $\operatorname{Wh}(C_5)$. To przykład
przeszkody, której nie może wytworzyć żadna macierz
nad samym $\Z$.
\end{przyklad}

\section{Torsja acyklicznego kompleksu i prosta homotopia}
\label{sec:torsja-kompleksu-20}

Kompleks względny kobordyzmu jest acykliczny nad $R$.
Acykliczność pozwala tu zbudować homotopię łańcuchową
od identyczności do zera, lecz samo jej stwierdzenie nie
mówi, czy odwracalne macierze da się zbudować z
ruchów elementarnych. Następna definicja zachowuje
tę różnicę.

\begin{lemat}[Kontrakcja skończonego kompleksu]
\label{lem:kontrakcja-20}
Niech $C_\bullet$ będzie skończonym acyklicznym
kompleksem skończenie generowanych wolnych modułów
nad pierścieniem $R$. Istnieją homomorfizmy
$s_i:C_i\to C_{i+1}$ takie, że
\[
 \partial s+s\partial=\id,
 \qquad s^2=0.
\]
\end{lemat}
\begin{proof}
Oznaczmy $B_i=\operatorname{im}\partial_{i+1}$.
Zaczynając od najniższego niezerowego stopnia,
acykliczność daje surjekcję $C_{i+1}\to B_i=C_i$;
rozdziela się ona, bo $C_i$ jest wolny.
Jej jądro $B_{i+1}$ jest zatem składnikiem prostym
$C_{i+1}$, więc modułem rzutowym. Indukcyjnie
każda surjekcja $C_{j+1}\to B_j$ rozdziela się,
a dostajemy rozkład
$C_j\cong B_j\oplus B_{j-1}$.
Różniczka przenosi drugi składnik identycznością
na $B_{j-1}$ w stopniu niższym. Definiujemy $s$
jako odwrotne przeniesienie pierwszego składnika
$B_j$ do drugiego składnika w stopniu wyższym;
na $B_{j-1}$ kładziemy $s=0$.
Na każdym z obu składników bezpośrednie obliczenie
daje $\partial s+s\partial=\id$ i $s^2=0$.
\end{proof}

\begin{definicja}[Torsja kompleksu]
\index{torsja kompleksu@Torsja kompleksu}
\label{def:torsja-kompleksu-20}
Dla acyklicznego kompleksu $C_\bullet$ z ustalonymi
bazami niech $s$ będzie kontrakcją z poprzedniego
lematu. Operator
\[
 T=\partial+s:C_{\mathrm{odd}}\longrightarrow
 C_{\mathrm{even}},\quad
 C_{\mathrm{odd}}=\bigoplus_{i\text{ nieparzyste}}C_i,
 \quad C_{\mathrm{even}}=\bigoplus_{i\text{ parzyste}}C_i
\]
jest izomorfizmem. Obie wolne sumy mają tę samą rangę:
dla $R=\Z[\pi]$ wynika to po zastosowaniu
$-\otimes_R\Z$, gdzie lewe działanie na $\Z$ określa augmentacja.
Izomorfizm przechodzi wtedy na izomorfizm wolnych grup abelowych,
które mają jednoznacznie określoną rangę.
\emph{Torsją} $C_\bullet$
nazywamy klasę $\tau(C_\bullet)=[T]\in K_1(R)$,
a w sytuacji geometrycznej jej obraz w
$\operatorname{Wh}(\pi)$. Ustalony porządek baz
w sumach prostych jest częścią konwencji; jego
zmiana dodaje najwyżej klasę $[(-1)]$, która znika
w grupie Whiteheada. Domyślnie sumujemy stopnie w kolejności rosnącej.
\end{definicja}

\begin{stwierdzenie}[Poprawność definicji i para]
\label{prop:torsja-poprawna-20}
Klasa $[T]\in K_1(R)$ nie zależy od wyboru
kontrakcji $s$. Dla kompleksu mającego jedynie
$C_{k+1}=R^r$, $C_k=R^r$ i różniczkę opisaną
odwracalną macierzą $A$ zachodzi
\begin{equation}\label{eq:torsja-dwa-stopnie-20}
 \tau(C_\bullet)=(-1)^k[A].
\end{equation}
W szczególności dla pary indeksów $2,3$ torsją
jest $[A]$.
\end{stwierdzenie}
\begin{proof}
Ponieważ $\partial^2=s^2=0$ i
$\partial s+s\partial=\id$, ten sam operator
$\partial+s$ w przeciwną stronę jest odwrotnością
$T$; definicja istotnie daje element $K_1(R)$.
Niech $s,t$ będą dwiema kontrakcjami o kwadracie zerowym.
Na całym kompleksie połóżmy
\[
 F=\partial s+t\partial=I+(t-s)\partial=I+N.
\]
Operator $N$ zachowuje stopień, ma obraz w
$B_i=\ker\partial_i$ i znika na $B_i$. Stąd $N^2=0$,
więc $F^{-1}=I-N$. Ponadto
\[
 F\partial=\partial F=\partial,
 \qquad Fs=tF=t\partial s.
\]
Zatem $F$ jest izomorfizmem kompleksów, który przeprowadza
kontrakcję $s$ w $t$.

Sprawdźmy, że w każdym stopniu $[F_i]=0$ w $K_1(R)$.
Operator $P_i=\partial_{i+1}s_i$ jest projekcją na $B_i$,
więc $P_iN_i=N_i$ i $N_iP_i=0$. W bazie wolnego modułu
$C_i$ są to zwykłe macierze, nawet jeśli sam $B_i$ nie jest wolny.
Mnożąc cztery macierze blokowe, otrzymujemy
\[
 \begin{pmatrix}I&P_i\\0&I\end{pmatrix}
 \begin{pmatrix}I&0\\N_i&I\end{pmatrix}
 \begin{pmatrix}I&-P_i\\0&I\end{pmatrix}
 \begin{pmatrix}I&0\\-N_i&I\end{pmatrix}
 =\begin{pmatrix}I+N_i&0\\0&I\end{pmatrix}.
\]
Każdy czynnik z lewej jest iloczynem macierzy elementarnych,
zatem $F_i$ jest elementarna po stabilizacji.
Z relacji $T_tF_{\mathrm{odd}}=F_{\mathrm{even}}T_s$ dostajemy
\[
 [T_t]+[F_{\mathrm{odd}}]=[F_{\mathrm{even}}]+[T_s].
\]
Obie klasy $F$ są sumami zerowych klas $[F_i]$.
Wynika stąd $[T_t]=[T_s]$, bez założenia wolności modułów brzegów.

Jeżeli występują tylko dwa stopnie, acykliczność
mówi, że $A$ jest izomorfizmem. Dla parzystego $k$
źródłem $T$ jest $C_{k+1}$ i $T=A$; dla nieparzystego
$k$ źródłem jest $C_k$ i $T=A^{-1}$.
Ponieważ $[A^{-1}]=-[A]$, dostajemy~\eqref{eq:torsja-dwa-stopnie-20}.
\end{proof}

\begin{uwaga}[To nie jest torsja elementu modułu]
W Definicji~\ref{def:torsja} element grupy abelowej,
czyli $\Z$-modułu, był torsyjny, gdy pewna niezerowa liczba
całkowita go zabijała. Tutaj \emph{torsja Whiteheada} jest klasą
odwracalnej macierzy całego acyklicznego kompleksu.
Może być niezerowa nawet wtedy, gdy wszystkie grupy
homologii są zerowe. Nazwy odnoszą się do różnych
konstrukcji.
\end{uwaga}

\begin{definicja}[Prosta równoważność homotopijna]
\index{prosta rownowaznosc homotopijna@Prosta równoważność homotopijna}
\label{def:prosta-homotopia-20}
\emph{Elementarne zwinięcie} skończonego kompleksu CW
usuwa dwie komórki wymiarów $k,k+1$, z których
pierwsza jest wolną ścianą drugiej; ruch odwrotny
to \emph{elementarne rozszerzenie}.
Precyzyjnie, przy takim zwinięciu $X$ do podkompleksu $Y$ mamy
$X=Y\cup e^k\cup e^{k+1}$, a odwzorowanie charakterystyczne
$\chi:D^{k+1}\to X$ przenosi wnętrze jednej zamkniętej
półkuli w $\partial D^{k+1}$ homeomorficznie na $e^k$
i cały pozostały brzeg do $Y$. Pozwala to zsunąć obie komórki
na $Y$; na przykład trójkąt zwija się przez usunięcie swojego
wnętrza i wnętrza wolnej krawędzi na dwie pozostałe krawędzie.
Równoważność
homotopijną $f:Y\to X$ między spójnymi skończonymi kompleksami CW
nazywamy \emph{prostą równoważnością homotopijną}, gdy
po odpowiednich aproksymacjach komórkowych jej
klasę można otrzymać skończonym ciągiem tych
ruchów i izomorfizmów komórkowych. Równoważnie,
$\tau(f)=0\in\operatorname{Wh}(\pi_1X)$, gdzie
$\tau(f)$ definiuje torsja acyklicznego stożka
odwzorowania podniesionych kompleksów łańcuchów.
Dla inkluzji $Y\hookrightarrow X$ przyjmujemy
zgodną konwencję $\tau(X,Y)=\tau(C_*(\widetilde X,
\widetilde Y))$.
\end{definicja}

\begin{uwaga}[Stożek w konwencji homologiczej]
Stożek zbiera w jednym kompleksie informację o tym, czy
odwzorowanie łańcuchowe $F:D_\bullet\to C_\bullet$
indukuje izomorfizm homologii. Jego moduły i operator brzegu to
\[
 \operatorname{Cone}(F)_i=C_i\oplus D_{i-1},\qquad
 \partial(c,d)=(\partial_Cc+F(d),-\partial_Dd).
\]
Równość $\partial^2=0$ wynika z $\partial_CF=F\partial_D$.
Ciąg dokładny z podkompleksem $C_\bullet$ i ilorazem
$D_{\bullet-1}$ pokazuje, że stożek jest acykliczny dokładnie
wtedy, gdy $F_*$ jest izomorfizmem w każdym stopniu.
W powyższej definicji bierzemy $F=C_*(\widetilde f)$,
utożsamiając grupy podstawowe przez $f_*$;
bazy są sumami wybranych baz komórkowych $X$ i $Y$.
Przykładowo stożek identyczności modułu $R$ w stopniu $0$
jest kompleksem $R\xrightarrow{1}R$ w stopniach $1,0$,
o torsji zerowej. Jest to homologiczna wersja konstrukcji
użytej dla kochainów w Sekcji~\ref{sec:twierdzenie-de-rhama}.
\end{uwaga}

\begin{uwaga}[Skąd ta równoważność?]
\label{uwaga:wh-charakterystyka-20}
W pojedynczym zwinięciu względny kompleks pary
komórek ma różniczkę $(\pm g)$: odpowiedni element
$\operatorname{Wh}(\pi)$ jest zerowy. Rachunek stożków
pokazuje addytywność torsji przy składaniu ruchów.
Głębsza część charakterystyki mówi odwrotnie, że
zerowa klasa torsji pozwala po stabilizacjach i
zmianach komórkowych rozłożyć równoważność na takie
ruchy. Jest to twierdzenie Whiteheada dla równoważności
spójnych skończonych kompleksów CW, przyjęte tutaj jako wynik
zewnętrzny; zob.
\href{https://webhomes.maths.ed.ac.uk/~v1ranick/papers/owsem.pdf}
{M.\ Kreck i W.\ Lück, \emph{The Novikov Conjecture}, twierdzenie 6.15(1)}.
Nie jest to definicja acykliczności. W naszym dowodzie
kobordyzmu zrealizujemy odpowiadającą mu redukcję
bezpośrednio uchwytami.
\end{uwaga}

\section{Torsja kobordyzmu i postać normalna uchwytów}
\label{sec:torsja-kobordyzmu-20}

\begin{definicja}[Torsja $h$-kobordyzmu]
\label{def:torsja-hc-20}
\index{torsja Whiteheada@Torsja Whiteheada}
Jeżeli $(W;N_-,N_+)$ jest $h$-kobordyzmem, to
\[
 \tau(W,N_-):=\tau\bigl(C_*(\widetilde W,
     \widetilde N_-;\Z)\bigr)
 \in\operatorname{Wh}(\pi_1N_-).
\]
Grupę $\pi_1N_-$ utożsamiamy z $\pi_1W$ przez
inkluzję. Mówimy o \emph{$s$-kobordyzmie}, gdy
ta klasa jest zerowa.
\end{definicja}

\begin{stwierdzenie}[Niezależność od rozkładu]
\label{prop:niezaleznosc-torsji-20}
Torsja $\tau(W,N_-)$ nie zależy od rozkładu na
uchwyty, triangulacji zgodnej ze strukturą gładką, orientacji rdzeni ani wyboru
podniesień. Jest niezmiennikiem dyfeomorfizmu
kobordyzmów względnie $N_-$.
\end{stwierdzenie}
\begin{proof}
Zmiana orientacji i podniesienia generatora
$e$ daje $e\mapsto e(\pm g)$, a klasa takiej
zmiany znika w $\operatorname{Wh}(\pi)$.
Przesunięcie uchwytu zmienia bazę macierzą
$E_{ij}(a)$, więc również nie zmienia klasy.
Narodziny lub zniesienie pary uchwytów po odpowiednich
elementarnych zmianach baz dołączają lub usuwają
acykliczny składnik $R\xrightarrow{\ \pm g\ }R$;
z~\eqref{eq:torsja-dwa-stopnie-20} jego torsja
znika w $\operatorname{Wh}(\pi)$.
Permutacja rdzeni dodaje klasę $[(-1)]$, także usuwaną w ilorazie.
Te rachunki sprawdzają poszczególne ruchy, lecz same nie
porównują dowolnych rozkładów.

Używamy więc jawnie następującego wyniku zewnętrznego:
dla zwartego gładkiego $h$-kobordyzmu torsja kompleksu
dowolnego uporządkowanego rozkładu Morse'a jest równa torsji
pary komórkowej z dowolnej triangulacji zgodnej ze strukturą
gładką i z brzegiem. Ta ostatnia klasa nie zależy od takiej
triangulacji. Dokładne porównanie daje
\href{https://webhomes.maths.ed.ac.uk/~v1ranick/papers/milnorwh.pdf}
{J.\ Milnor, \emph{Whitehead Torsion}, lemat 9.1 i twierdzenie 9.3}.
Ich wejściem geometrycznym jest teoria gładkich triangulacji,
a algebraicznym niezmienniczość torsji przy podziale komórek.
Wynik ten, razem z powyższymi rachunkami baz, daje niezależność
od wszystkich wymienionych wyborów. Dyfeomorfizm przenosi
triangulację i jej podniesienia, co dowodzi niezmienniczości
względem $N_-$.
\end{proof}

\begin{lemat}[Postać normalna $h$-kobordyzmu]
\label{lem:postac-normalna-20}
Przy $m\geq6$ rozkład uchwytowy $h$-kobordyzmu
względem $N_-$ można zmienić, zachowując sam
kobordyzm względem $N_-$, tak aby pozostały tylko
uchwyty indeksów $2$ i $3$. Dokładniej, dla każdego
$2\leq k\leq m-3$ można otrzymać rozkład zawierający
tylko indeksy $k,k+1$.
\end{lemat}
\begin{proof}
Porządkujemy indeksy i usuwamy $0$-uchwyty przy
spójnym brzegu wejściowym, jak w
Lemacie~\ref{lem:zero-m-hc-19}. Odwrócony rozkład
usuwa $m$-uchwyty przy $N_+$. Zamiana $1$-uchwytu
na $3$-uchwyt z
Lematu~\ref{lem:handel-trade-hc-19} wymaga teraz
uważniejszego doboru okręgu. Zaczynamy od pętli
przechodzącej raz przez sferę pasa wybranego
$1$-uchwytu. Jej obraz w $\pi_1(W)$ pochodzi,
na mocy surjektywności $\pi_1(N_-)\to\pi_1(W)$,
od pewnej pętli w $N_-$. Do łuku zamykającego
doklejamy reprezentanta odwrotności tej klasy
(z odpowiednią drogą łączącą punkty bazowe).
Nadal przecina ona wybraną sferę pasa raz,
lecz jej klasa w $\pi_1(W)$ jest teraz zerowa.
Po przeniesieniu za stare $2$-uchwyty okrąg
pozostaje homotopijnie trywialny w poziomicy:
późniejsze uchwyty indeksu co najmniej $3$
nie zmieniają $\pi_1$, a dualny opis dolnej
części też nie zmienia $\pi_1$ tej poziomicy.
Przeniesiony okrąg ogranicza zatem zanurzony dysk:
w poziomicy wymiaru co najmniej $5$ względna pozycja ogólna
usuwa samoprzecięcia dwuwymiarowego wypełnienia.
Koncentryczne okręgi tego dysku dają izotopię do małego okręgu.
Tam tworzymy parę $2,3$ ze standardowym obramowaniem;
przenosimy je razem z izotopią i przepływem w dół.
Jak w dowodzie Lematu~\ref{lem:handel-trade-hc-19}, nowe
przyklejenie omija sfery pasa starych $2$-uchwytów.
Możemy więc ustawić nowy $2$-uchwyt tuż po wybranym
$1$-uchwycie i znieść je przez Twierdzenie~\ref{thm:zniesienie-18}.
Analogiczna operacja od $N_+$ usuwa indeks $m-1$.

Aby usunąć kolejny najmniejszy indeks $j\geq2$,
pracujemy po jego uchwytach, w poziomicy $L_j$.
Nie ma już generatorów $C_{j-1}$, a acykliczność
kompleksu nad $R$ daje
$\operatorname{im}\partial_{j+1}=C_j$.
Dla wybranego $j$-uchwytu o generatorze $e_1$
istnieją więc współczynniki $x_h\in R$ takie, że
\[
 e_1=\partial_{j+1}\Bigl(\sum_h e_h^{j+1}x_h\Bigr).
\]
Wybieramy w $L_j$ małą obramowaną sferę $S^j$,
która ogranicza dysk i nie spotyka sfer pasa.
Na poziomie po starych $(j+1)$-uchwytach
tworzymy lokalną znoszącą się parę
$(j+1,j+2)$ z taką sferą przyczepienia.
Teraz modyfikujemy tę sferę \emph{niżej},
w $L_j$: dla każdego jednomianu $\pm g$
we współczynniku $x_h$ wykonujemy sumę
spójną przez wąską rurę z równoległą kopią
sfery przyczepienia starego $(j+1)$-uchwytu
$h$. Orientacja kopii daje znak, a droga
prowadzenia rury wyznacza $g$.
Po dołączeniu starego uchwytu $D^{j+1}\times D^{m-j-1}$
ta równoległa kopia ogranicza dysk
$D^{j+1}\times\{v\}$, gdzie $v\in S^{m-j-2}$.
Jest to dysk \emph{w nowym brzegu}, a nie rdzeń leżący
we wnętrzu kobordyzmu. Przesunięcie rury przez ten dysk
cofa sumę spójną w wyższej poziomicy.
Dlatego zmodyfikowana sfera jest tam izotopijna do małej
sfery; razem z ruchem przenosimy obramowanie nowo narodzonej
pary. W dolnej poziomicy
jej klasa przecięć zmienia się dokładnie o
$\partial_{j+1}(e_h^{j+1}x_h)$.
Po skończonej liczbie rur klasa ma więc
postać $e_1$. Pozycja ogólna pozwala wybrać
rozłączne rury; ich obramowania przedłużają
obramowania kopiowanych sfer. Zadane etykiety można uzyskać,
ponieważ $\pi_1(L_j)\to\pi_1(W)$ jest izomorfizmem
(argument w następnym lemacie), a usunięcie sfer
przyczepienia o kowymiarze co najmniej $3$ nie zmienia $\pi_1$.
W tej indukcji mamy $j\leq m-4$, więc ten warunek zachodzi.
Rury można też odsunąć od sfer pasa: jednowymiarowy łuk
omija podrozmaitości kowymiaru co najmniej $2$.
Trik Whitneya z kontrolą grupy, opisany
w następnym lemacie, usuwa przeciwne składniki
o tej samej etykiecie. Sfera spotyka wtedy
sferę pasa wybranego $j$-uchwytu raz i omija
pozostałe. Nowy $(j+1)$-uchwyt ustawiamy przed starymi
uchwytami tego indeksu: jego opis w wyższej poziomicy
pochodzi z przepływu przyklejenia omijającego ich sfery pasa.
Wybrany $j$-uchwyt ustawiamy jako ostatni w swoim bloku.
Znosimy te dwa sąsiadujące uchwyty;
nowy $(j+2)$-uchwyt zostaje.
Powtarzamy dla wszystkich starych
$j$-uchwytów. Ten argument jest geometryczną
realizacją surjektywności różniczki, a nie
samym wnioskowaniem z równości grup homologii.

Indukcja od dołu usuwa indeksy poniżej wybranego
$k$. Następnie odwracamy kobordyzm i powtarzamy
argument od $N_+$, usuwając indeksy powyżej
$k+1$. Przy odwróceniu $j$-uchwyt ma indeks
$m-j$. Usuwanie pierwotnego indeksu $l\geq k+2$
od góry może dodać indeks $l-2\geq k$, więc nie odtwarza
już usuniętych indeksów poniżej $k$.
Warunek $2\leq k\leq m-3$ zapewnia
potrzebną przestrzeń do izotopii i dysków
Whitneya. Ostatecznie pozostają tylko
indeksy $k,k+1$. W przypadku skrajnym $k=2$
używamy opisu dopełnienia z dowodu
następnego lematu; warunek obramowania
sprawdzamy za pomocą geometrycznej wersji
triku Whitneya omówionej w
Rozdziale~\ref{ch:hcobordyzm}.
\end{proof}

W postaci normalnej $2,3$ acykliczny kompleks ma
jedyną różniczkę
\[
 0\longrightarrow R^r\xrightarrow{\ A\ }R^r
 \longrightarrow0.
\]
Różniczka jest izomorfizmem przez acykliczność.
Zastosowanie $-\otimes_R\Z$ przez augmentację pokazuje,
że liczby uchwytów są równe, a więc $A\in GL_r(R)$.
Z~\eqref{eq:torsja-dwa-stopnie-20} dostajemy
\begin{equation}\label{eq:torsja-macierz-20}
             \tau(W,N_-)=[A]\in\operatorname{Wh}(\pi).
\end{equation}
To jest najprostsza postać torsji: jedna odwracalna
macierz, której klasy nie da się odczytać z
$H_*(W,N_-;\Z)$.

\begin{lemat}[Trik Whitneya z etykietami]
\label{lem:whitney-pi-20}
Niech $m\geq6$, a uporządkowany rozkład $h$-kobordyzmu
ma wszystkie indeksy w przedziale $[k,m-2]$,
gdzie $2\leq k\leq m-3$. Niech $L$ będzie poziomicą
po wszystkich $k$-uchwytach, przed wszystkimi
$(k+1)$-uchwytami. Dwa przecięcia sfer
$A_j$ i $B_i$ o znakach przeciwnych i tej samej
etykiecie $g\in\pi$ można usunąć izotopią,
nie tworząc nowych przecięć z pozostałymi sferami.
Jeżeli $a_{ij}=\pm g$, po skończonej liczbie
powtórzeń można otrzymać dokładnie jeden
punkt przecięcia tej pary sfer. Dla $a_{ij}=0$ można
usunąć wszystkie ich przecięcia. Założenia nie wymagają
postaci normalnej: dopuszczają również wyższe indeksy.
\end{lemat}
\begin{proof}
Najpierw $\pi_1(L)\to\pi_1(W)$ jest izomorfizmem.
Dolną podpoziomicę otrzymujemy od $L$ przez uchwyty dualne
indeksu co najmniej $m-k\geq3$, a od niej do $W$
dołączamy uchwyty indeksu co najmniej $k+1\geq3$.
Żaden z tych kroków nie zmienia grupy podstawowej.

Wybieramy łuk między punktami na $A_j$ i łuk
powrotny na $B_i$, unikając innych przecięć.
Definicja etykiety w~\eqref{eq:macierz-pi-20}
pokazuje, że klasa ich złożenia jest
$g_xg_y^{-1}$, z dokładnością do sprzężenia i wyboru
kierunku obchodzenia. Ponieważ $g_x=g_y$,
pętla jest homotopijnie trywialna. Kluczowe
jest znalezienie jej wypełnienia w
\emph{dopełnieniu} całej rodziny sfer.
Oznaczmy to dopełnienie przez
$U=L\setminus(\bigcup A_j\cup\bigcup B_i)$.
Gdy $2<k<m-3$, obie rodziny mają kowymiar co najmniej $3$.
Pozycja ogólna pozwala odsunąć od nich zarówno drogi,
jak i dwuwymiarowe homotopie dróg. Stąd
$\pi_1(U)\to\pi_1(L)$ jest izomorfizmem.

Dla $k=2$ dolna część zawiera tylko $2$-uchwyty.
Niech $E$ będzie dopełnieniem wnętrz ich obramowanych
otoczeń przyczepienia $S^1\times D^{m-2}$ w $N_-$.
Zarówno $N_-$ bez okręgów przyczepienia, jak i
$L\setminus\bigcup B_i$ retrahują się na $E$:
po stronie nowych uchwytów kurczymy promieniowo
$(D^2\setminus\{0\})\times S^{m-3}$ ku
$S^1\times S^{m-3}$. Usunięcie okręgów o kowymiarze
$m-2\geq4$ zachowuje $\pi_1(N_-)$.
Inkluzje $E\to N_-\to W$ dają więc izomorfizm,
a przez powyższe retrakcje również
$\pi_1(L\setminus\bigcup B_i)\to\pi_1(W)$
jest izomorfizmem. Sfery $A_j$ mają kowymiar
$m-3\geq3$, więc ich usunięcie zachowuje ten wynik.
Dla $k=m-3$ górna część zawiera tylko uchwyty indeksu
$m-2$. Odwracamy ją: są to $2$-uchwyty, których
sferami pasa w $L$ są właśnie $A_j$. Ten sam argument
od $N_+$ daje izomorfizm dla $L\setminus\bigcup A_j$,
a potem dla $U$. Wszystkie te homomorfizmy są indukowane
przez inkluzje do tego samego $W$; nie porównujemy jedynie
abstrakcyjnie izomorficznych grup.

Odsunięta od sfer kopia pętli Whitneya jest zatem ściągalna
w $U$. Wypełniamy ją odwzorowaniem dysku, a względna pozycja
ogólna daje zanurzenie: $2+2-(m-1)<0$.
Dołączamy kołnierz ku łukom na sferach, zachowując standardowe
naroża w obu przecięciach. Otrzymujemy dysk o czystym wnętrzu.

Obramowanie i sam ruch zapewnia dokładnie zewnętrzny lemat
Milnora przytoczony w Kroku 3 dowodu
Lematu~\ref{lem:whitney-hc-19}.
Dla $k\leq m-4$ stosujemy go do
$P=\bigcup A_j$, $Q=\bigcup B_i$, o wymiarach
$r=k\geq2$ i $s=m-k-1\geq3$.
Przy $r=2$ sprawdziliśmy wymaganą injektywność
$\pi_1(L\setminus Q)\to\pi_1(L)$.
Dla $k=m-3$ zamieniamy role rodzin, a na końcu
odwracamy izotopię, aby poruszać $A$ względem $B$.
Orientacje sfer i ich wiązek normalnych pochodzą z uchwytów.
Równość etykiet zapewnia zgodność lokalnych znaków wzdłuż
ściągalnej pętli, także gdy $W$ jest nieorientowalna.
Zastosowanie lematu do całych rodzin chroni pozostałe przecięcia.
Nie wywodzimy obramowania jedynie z trywialności wiązki nad dyskiem.

Współczynnik $a_{ij}=\pm g$ jest sumą skończoną.
Po pogrupowaniu monomów każdy nadmiarowy
składnik ma partnera $\pm g\mp g$; powtarzamy
poprzedni ruch, aż zostaje jeden punkt
z etykietą $g$ i właściwym znakiem.
\end{proof}

\section{Twierdzenie o \texorpdfstring{$s$}{s}-kobordyzmie i jego dowód}
\label{sec:tw-s-kobordyzm-20}

\begin{twierdzenie}[Mazur--Stallings--Barden]
\label{thm:s-kobordyzm-20}
Niech $m\geq6$ i $(W^m;N_-,N_+)$ będzie gładkim
$h$-kobordyzmem. Wtedy następujące warunki są równoważne:
\begin{enumerate}[label=(\roman*)]
 \item $\tau(W,N_-)=0\in\operatorname{Wh}(\pi_1N_-)$;
 \item istnieje dyfeomorfizm
 $\Phi:N_-\times[0,1]\to W$ z $\Phi(x,0)=x$.
\end{enumerate}
\end{twierdzenie}

\begin{proof}
\emph{Krok 1: z dowolnego rozkładu do macierzy.}
Lemat~\ref{lem:postac-normalna-20} daje rozkład
z samymi uchwytami indeksów $2,3$.
Ponieważ względny kompleks nakrycia jest
acykliczny, jego macierz $A$ jest odwracalna.
Wzór~\eqref{eq:torsja-macierz-20} utożsamia
założenie (i) z $[A]=0$ w $\operatorname{Wh}(\pi)$.

\emph{Krok 2: redukcja stabilna.}
Z definicji grupy Whiteheada $[A]=0$ oznacza,
że po dołączeniu bloku jednostkowego $I_s$
można sprowadzić $A\oplus I_s$ do $I_{r+s}$
generowanymi ruchami: dodawaniem do wiersza
lub kolumny innego wiersza lub kolumny
pomnożonej przez element $R$ (z lewej dla wierszy,
z prawej dla kolumn), zamianą kolejności
oraz mnożeniem wybranej bazy przez $\pm g$.
Dlaczego skończona sekwencja wystarcza?
W $GL_\infty(R)$ równość klas oznacza równość
po skończonej liczbie stabilizacji i przemnożeń
przez macierze elementarne; iloraz przez
$[\pm g]$ dodaje tylko skończenie wiele zmian
pojedynczych baz. Przykład
$[(u)]\ne0$ z~\ref{prz:wh-c5-20} pokazuje,
że sam fakt odwracalności $A$ nie daje tej redukcji.

\emph{Krok 3: realizacja ruchów.}
Blok $I_s$ tworzymy $s$ razy przez narodziny
znoszącej się pary uchwytów $2,3$. Zmiana bazy
źródła $C_3$ przez przesunięcie $3$-uchwytów
mnoży macierz różniczki z prawej przez macierz
elementarną; przesunięcia $2$-uchwytów zmieniają
bazę celu $C_2$ i odpowiadają mnożeniu z lewej
przez odwrotną macierz elementarną. Dokładniej, zapisując
bazy jako wiersze generatorów, dla $e'=eU$, $f'=fV$ mamy
\[
 \partial(e')=\partial(e)U=fAU=f'V^{-1}AU,
 \qquad A'=V^{-1}AU.
\]
Pas łączący
uchwyty prowadzimy wzdłuż pętli reprezentującej
$g$, aby uzyskać współczynnik $\pm g$ w tej konwencji.
Taką drogę można wybrać, bo odpowiednia poziomica ma tę
samą grupę podstawową co $W$. Dla $3$-uchwytów sfery
przyczepienia mają kowymiar co najmniej $3$;
dla $2$-uchwytów są okręgami kowymiaru co najmniej $4$.
Możemy więc odsunąć drogę od pozostałych przyklejeń bez
zmiany jej klasy. Obramowaną rurę wybieramy jak w
Stwierdzeniu~\ref{prop:slide-18}, przenosząc ramki wzdłuż łuku.
Element
$a=\sum n_g g\in R$ otrzymujemy skończoną
liczbą takich przesunięć. Właściwy porządek
mnożenia współczynników zachowujemy przez
konwencję prawych modułów z
Sekcji~\ref{sec:pi-lancuchy-20}. Permutacja
uchwytów i odwrócenie orientacji zmieniają
bazę, a zmiana podniesienia realizuje czynnik
$g$. Ruchy te pozostawiają dyfeomorfizm
$W$ względnie $N_-$.
Na koniec otrzymujemy macierz $I_{r+s}$.

\emph{Krok 4: z jedynek do punktów.}
Równość macierzy jednostkowej oznacza tylko,
że $A_j$ i $B_i$ mają sumy etykietowane
$\delta_{ij}e$. Dla wszystkich par $i\ne j$
Lemat~\ref{lem:whitney-pi-20} usuwa
wszystkie ich przecięcia parami o znaku
przeciwnym i tej samej etykiecie.
Dla $i=j$ pozostaje dokładnie jeden punkt
z etykietą $e$. Wewnętrzne dyski Whitneya
można dobierać kolejno bez wprowadzania
nowych przecięć, więc proces kończy się.
Twierdzenie~\ref{thm:zniesienie-18}
usuwa teraz kolejno wszystkie geometryczne
pary $2,3$. Funkcja Morse'a bez punktów
krytycznych daje wymagany dyfeomorfizm
na mocy Lematu~\ref{lem:brak-krytycznych-hc-19}.
To dowodzi (i)$\Rightarrow$(ii).

\emph{Krok 5: konieczność.}
Jeśli (ii) zachodzi, przenosimy funkcję
$(x,t)\mapsto t$ z cylindra do $W$.
Nie ma ona punktów krytycznych i jej
względny kompleks uchwytów jest zerowy,
a więc ma torsję $0$. Z niezależności
torsji od rozkładu, udowodnionej w
Stwierdzeniu~\ref{prop:niezaleznosc-torsji-20},
wynika $\tau(W,N_-)=0$.
\end{proof}

\begin{figure}[htbp]
\centering
\begin{tikzpicture}[>=Latex,font=\small,
 every node/.style={align=center,rounded corners=3pt,inner sep=5pt}]
 \node[draw=manifoldblue!75!black,fill=manifoldblue!9]
 (a) at (-3.15,.90) {$h$-kobordyzm\\$C_*(\widetilde W,\widetilde N_-)$};
 \node[draw=manifoldgold!80!black,fill=manifoldgold!12]
 (b) at (2.65,.90) {$2$- i $3$-uchwyty\\$A\in GL_r(\Z[\pi])$};
 \node[draw=manifoldgreen!75!black,fill=manifoldgreen!10]
 (c) at (2.65,-.90) {jeśli $[A]=0$\\stabilizacja i przesunięcia};
 \node[draw=manifoldred!75!black,fill=manifoldred!9]
 (d) at (-3.15,-.90) {trik Whitneya i zniesienia\\$W\cong N_-\times[0,1]$};
 \draw[->,thick] (a)--(b);
 \draw[->,thick] (b)--(c);
 \draw[->,thick] (c)--(d);
\end{tikzpicture}
\caption{Logika dowodu twierdzenia o $s$-kobordyzmie.
Acykliczność daje odwracalność $A$; znikanie jej klasy
w grupie Whiteheada daje dopuszczalną redukcję do $I$.
Na końcu potrzebny jest jeszcze geometryczny trik Whitneya.}
\label{fig:dowod-s-kobordyzm-20}
\end{figure}

\begin{wniosek}[Odzyskanie twierdzenia Smale'a]
Jeżeli $\pi_1(W)=\{e\}$, to
$\operatorname{Wh}(\pi_1W)=0$ na mocy
Stwierdzenia~\ref{prop:wh-trywialna-20}.
Twierdzenie~\ref{thm:s-kobordyzm-20}
daje zatem Twierdzenie~\ref{thm:smale-hc-19}.
\end{wniosek}

\begin{uwaga}[Który brzeg i orientacja?]
Torsję zdefiniowaliśmy względem $N_-$.
Względem $N_+$ istnieje analogiczna torsja;
łączy je twierdzenie o dualności (ze znakiem
i inwolucją $g\mapsto g^{-1}$, skorygowaną
przez charakter orientacji w przypadku
nieorientowalnym). W szczególności
\emph{zerowanie} jednej jest równoważne
zerowaniu drugiej. W dowodzie potrzebowaliśmy
tylko torsji względem wybranego brzegu wejściowego,
więc formuła znaków dualności nie jest potrzebna.
\end{uwaga}

\section{Przykłady i miejsce w teorii chirurgii}
\label{sec:przyklady-s-kobordyzm-20}

\begin{przyklad}[Para pozornie skomplikowanych uchwytów]
Rozważmy kompleks z $2$- i $3$-uchwytami,
którego macierz nad $R$ jest
\[
 A=\begin{pmatrix}1&t\\0&1\end{pmatrix},\qquad t\in\pi.
\]
Jej odwrotność istnieje, ale ważniejsze jest
$A=E_{12}(t)$. Klasa $[A]$ jest zerowa bez
stabilizacji. Jedno przesunięcie uchwytu
usuwa wyraz $t$, potem kontrolowany trik
Whitneya sprowadza przecięcia do dwóch
par po jednym punkcie. Uchwyty można znieść.
\end{przyklad}

\begin{przyklad}[Niezerowa torsja mimo zerowej homologii]
W przykładzie $C_5$ z~\ref{prz:wh-c5-20}
weźmy abstrakcyjny kompleks w stopniach $3,2$
\[
 0\longrightarrow R\xrightarrow{\ u\ }R
   \longrightarrow0,\qquad u=1-t-t^{-1}.
\]
Jest acykliczny, bo $u$ ma odwrotność $v$.
Ma jednak torsję $[(u)]\ne0$.
Po augmentacji różniczka jest mnożeniem
przez $\epsilon(u)=-1$, więc również zwykły
kompleks całkowity jest acykliczny. Ten przykład
oddziela trzy własności: znikanie zwykłej
homologii, acykliczność nad $R$ oraz znikanie
torsji. Dopiero ostatnia pozwala zastosować
twierdzenie o $s$-kobordyzmie.
\end{przyklad}

\begin{lemat}[Addytywność przy sklejaniu]
\label{lem:add-torsja-20}
Jeśli $(W;N_0,N_1)$ oraz $(V;N_1,N_2)$ są
$h$-kobordyzmami, to po przeniesieniu obu klas
do $\operatorname{Wh}(\pi_1(N_0))$ zachodzi
\[
 \tau(W\cup_{N_1}V,N_0)
 =\tau(W,N_0)+\tau(V,N_1).
\]
\end{lemat}
\begin{proof}
Dobieramy zgodne struktury komórkowe na wspólnym
brzegu oraz podnosimy je do tego samego nakrycia
uniwersalnego sklejenia. Filtracja par
$(W\cup V,N_0)\supset(W,N_0)$ daje krótki ciąg
kompleksów
\[
 0\to C_*(\widetilde W,\widetilde N_0)
 \to C_*(\widetilde{W\cup V},\widetilde N_0)
 \to C_*(\widetilde V,\widetilde N_1)\to0.
\]
Wybieramy najpierw bazy komórek z $W$, potem
z $V$. Na poziomie modułów ten ciąg rozdziela się.
W takim rozszczepieniu różniczka kompleksu
środkowego ma postać
$\partial_B=\left(\begin{smallmatrix}\partial_A&u\\0&\partial_C\end{smallmatrix}\right)$,
gdzie $\partial_Au+u\partial_C=0$. Niech $s_A,s_C$ będą
kontrakcjami skrajnych kompleksów o kwadracie
zerowym. Wybieramy
\[
 s_B=\begin{pmatrix}
 s_A&-s_Au s_C\\0&s_C
 \end{pmatrix}.
\]
Bezpośrednie mnożenie z użyciem
$\partial_As_A+s_A\partial_A=1$ i
$\partial_Cs_C+s_C\partial_C=1$
daje $\partial_Bs_B+s_B\partial_B=1$ oraz $s_B^2=0$.
Na przykład blok pozadiagonalny pierwszej równości wynosi
\[
 -\partial_As_Aus_C+us_C+s_Au-s_Aus_C\partial_C
 =s_Au(1-\partial_Cs_C-s_C\partial_C)=0.
\]
W drugim rachunku ten blok jest równy
$-s_A^2us_C-s_Aus_C^2=0$.
Operator $\partial_B+s_B$ między sumami stopni
nieparzystych i parzystych jest blokowo
trójkątny, z $\partial_A+s_A$ i $\partial_C+s_C$ na przekątnej,
gdy zgrupujemy najpierw wszystkie składniki $A$, potem $C$.
Blok pozadiagonalny usuwamy macierzami
elementarnymi, więc w
$K_1(R)$ klasa środkowa jest sumą dwóch
pozostałych w tym porządku baz. Powrót do porządku według
stopni może dodać klasę $[(-1)]$ za permutację, która znika
w $\operatorname{Wh}(\pi)$. Tam otrzymujemy żądany wzór.
Utożsamienie grup jest dane
przez inkluzje brzegów, będące równoważnościami
homotopijnymi.
\end{proof}

\begin{twierdzenie}[Realizacja i klasyfikacja $h$-kobordyzmów]
\label{thm:klasyfikacja-hc-20}
Jeśli $N$ jest spójną zamkniętą gładką rozmaitością
wymiaru co najmniej $5$, to każdy element
$\operatorname{Wh}(\pi_1N)$ jest torsją pewnego
gładkiego $h$-kobordyzmu wychodzącego z $N$.
Co więcej, torsja daje bijekcję między klasami
takich $h$-kobordyzmów modulo dyfeomorfizmy równe identyczności
na oznaczonym brzegu wejściowym $N$ a
$\operatorname{Wh}(\pi_1N)$.
\end{twierdzenie}
\begin{proof}
\emph{Realizacja.} Przedstawmy zadaną klasę przez
$A\in GL_r(R)$, gdzie $m=\dim N+1\geq6$.
Do $N\times[0,1]$ przyczepiamy $r$ uchwytów indeksu $2$
w rozłącznych kulach górnego brzegu. Ich okręgi są małymi
standardowymi okręgami, z obramowaniami pochodzącymi od dysków
w tych kulach. Ten wybór jest istotny: sama ściągalność okręgu
nie ustala obramowania. Oznaczmy powstałą dolną część przez
$W_2$, a jej górny brzeg przez $L$.
Lokalny model operacji zastępuje w brzegu kulę fragmentem
$S^2\times S^{m-3}$ z usuniętą kulą. Mamy więc bazowe sfery
$S^2\times\{z\}\subset L$ o trywialnych wiązkach normalnych.
Po wyborze dróg bazowych ich klasy $\mu_i\in\pi_2(L)$
przechodzą na bazę
\[
 H_2(\widetilde W_2,\widetilde N;\Z)\cong R^r.
\]
Widać to także bezpośrednio z kompleksu względnego:
ma on tylko $r$ generatorów w stopniu $2$.

Przez sumy klas $\mu_i$, zmianę orientacji i działanie
dróg reprezentujących $g\in\pi$ uzyskujemy dla każdej kolumny
$A$ odwzorowanie $S^2\to L$ o żądanej klasie względnej.
Tu $\pi_1(L)\to\pi_1(W_2)$ jest izomorfizmem, ponieważ
dualny rozkład ma indeks $m-2\geq4$.
Pozycja ogólna pozwala zastąpić wszystkie te odwzorowania
rozłącznymi zanurzeniami $S_j\subset L$ bez zmiany ich klas:
$2+2<\dim L$. Pozostaje sprawdzić obramowania.

Funkcja
\[
 \pi_2(L)\longrightarrow\Z/2,\qquad
 [f]\longmapsto\langle w_2(TL),f_*[S^2]\rangle
\]
jest addytywna i nie zmienia się pod działaniem $\pi_1(L)$.
Wynika to z opisu sumy sfer przez ściągnięcie równika
oraz z tego, że zmiana drogi bazowej nie zmienia niezorientowanej
modulo $2$ klasy homologycznej sfery w $L$.
Na każdej $\mu_i$ funkcja ta jest zerowa:
$TL|_{S^2}=TS^2\oplus\underline{\R}^{m-3}$,
a $w_2(TS^2)=e(TS^2)\bmod2=2\bmod2=0$.
Dlatego znika także na każdej skonstruowanej $S_j$.
Ze wzoru Whitneya dla $TL|_{S_j}=TS_j\oplus\nu_j$
otrzymujemy $w_2(\nu_j)=0$. Ponadto $w_1(\nu_j)=0$,
bo $H^1(S^2;\Z/2)=0$.

Wyjaśnijmy, dlaczego to wystarcza. Zorientowaną wiązkę rangi
$q\geq3$ nad $S^2$ klasyfikuje funkcja sklejająca w
$\pi_1(SO(q))$ (Sekcja~\ref{sec:clutching-16}).
Dla $q=3$ ta grupa jest $\Z/2$, a klasę wykrywa $w_2$,
jak w tamtejszym przykładzie. Wiązka
$SO(q)\to SO(q+1)\to S^q$ i jej ciąg homotopii pokazują,
że $\pi_1(SO(q))\to\pi_1(SO(q+1))$ jest izomorfizmem
dla $q\geq3$, gdyż $\pi_2(S^q)=\pi_1(S^q)=0$.
Klasa $w_2$ jest stabilna, więc to samo kryterium działa
w każdej randze $q\geq3$.
Nasze $\nu_j$ mają rangę $m-3\geq3$ i są zatem trywialne.
Wybieramy ich obramowania i przyczepiamy $r$ uchwytów indeksu $3$.
Otrzymany względny kompleks ma różniczkę $A$,
zatem jest acykliczny.
Uchwyty indeksów $2,3$ nie zmieniają
$\pi_1(N)$: okręgi przyczepienia $2$-uchwytów
są homotopijnie trywialne, a wyższe indeksy
nie działają na $\pi_1$. Podniesienia inkluzji do nakryć
uniwersalnych są więc odwzorowaniami przestrzeni po prostu
spójnych, a acykliczność kompleksu względnego daje izomorfizmy
homologii. Wniosek~\ref{cor:whitehead-homologiczny}
czyni podniesioną inkluzję równoważnością homotopijną.
Nakrycia zachowują wyższe grupy homotopii, a $\pi_1$ już
sprawdziliśmy; Twierdzenie~\ref{thm:whitehead} daje równoważność
$N\hookrightarrow W$. Po odwróceniu uchwytów
ich indeksy wynoszą $m-2,m-3\geq3$.
Dualna różniczka ma, z dokładnością do znaków baz,
macierz $A^*=(\overline{a_{ji}})$, gdzie
$\overline g=w(g)g^{-1}$, a $w(g)\in\{\pm1\}$ opisuje
zachowanie orientacji przy obchodzeniu pętli $g$.
Jest to lemat dualności kompleksów uchwytowych:
odwracamy drogę przecięcia i uwzględniamy transport
orientacji (\href{https://webhomes.maths.ed.ac.uk/~v1ranick/papers/milnorwh.pdf}
{Milnor, \emph{Whitehead Torsion}, \S10, w tym uwaga
o przypadku nieorientowalnym}). Macierz ta jest odwracalna,
z odwrotnością $(A^{-1})^*$, co daje
acykliczność po stronie $N_+$; inkluzja
$N_+\hookrightarrow W$ jest izomorfizmem
na $\pi_1$, więc to samo twierdzenie Whiteheada
czyni ją równoważnością. Powstał
$h$-kobordyzm, a~\eqref{eq:torsja-macierz-20}
daje mu torsję $[A]$.

\emph{Jednoznaczność względem wejścia.}
Niech $W_1:N\to N_1$ i $W_2:N\to N_2$ mają tę
samą torsję $\tau_0$. Dzięki części o realizacji
budujemy $V:N_1\to N'$ o torsji $-\tau_0$,
przeniesionej do $\operatorname{Wh}(\pi_1N_1)$
przez izomorfizm indukowany przez $W_1$.
Lemat~\ref{lem:add-torsja-20} daje
$\tau(W_1\cup_{N_1}V,N)=0$, więc
Twierdzenie~\ref{thm:s-kobordyzm-20}
identyfikuje ten złożony kobordyzm z
$N\times[0,1]$ względem $N$.
Na drugim brzegu otrzymujemy dyfeomorfizm
$\phi:N'\to N$ i za jego pomocą sklejamy
trzy kobordyzmy w kolejności
\[
 H=W_1\cup_{N_1}V\cup_\phi W_2.
\]
Ponieważ $W_1\cup V$ jest iloczynem,
$H$ jest dyfeomorficzny względem $N$
z $W_2$. Z drugiej strony przeniesienie
grup podstawowych przez ten iloczyn
sprawia, że $\tau(V\cup_\phi W_2,N_1)
=-\tau_0+\tau_0=0$.
Ta druga część $H$ jest więc iloczynem
względem $N_1$, a $H$ jest także
dyfeomorficzny względem $N$ z $W_1$.
Składając oba dyfeomorfizmy uzyskujemy
$W_1\cong W_2$ względem $N$.
\end{proof}

\begin{uwaga}[Przykład z $C_5$ staje się geometryczny]
W szczególności, jeśli $N$ ma
$\pi_1(N)\cong C_5$ i $\dim N\geq5$, element
$[(1-t-t^{-1})]\ne0$ z
Przykładu~\ref{prz:wh-c5-20} jest torsją
pewnego $h$-kobordyzmu o brzegu wejściowym
$N$. Ten kobordyzm nie jest iloczynem
względem $N$, mimo że oba brzegi mają ten
sam typ homotopijny co całe $W$.
\end{uwaga}

\begin{uwaga}[Granice wymiarowe i następny krok]
Dysk Whitneya w poziomicy wymiaru $4$ może
mieć samoprzecięcia, więc przedstawiony dowód
wymaga $m\geq6$. Gdy planujemy chirurgię
na rozmaitościach o zadanej grupie podstawowej,
$\tau\in\operatorname{Wh}(\pi)$ jest przeszkodą
w zamianie uzyskanego $h$-kobordyzmu na produkt.
W dalszej teorii pojawią się także przeszkody
chirurgiczne w grupach $L_*(\Z[\pi])$;
są to dodatkowe obiekty, określające możliwość
\emph{wytworzenia} odpowiedniego kobordyzmu.
Nie utożsamiamy ich z torsją Whiteheada.
\end{uwaga}

\par\smallskip
{\footnotesize\noindent\emph{Dalsza lektura:}
\href{https://webhomes.maths.ed.ac.uk/~v1ranick/papers/milnorwh.pdf}
{J.~Milnor, \emph{Whitehead Torsion}},
\href{https://qkhan.pages.iu.edu/research/SLIDES/Lecture02_KHAN.pdf}
{Q.~Khan, \emph{The $s/h$-Cobordism Theorem}}.\par}

\clearpage
